% U9 WS2 -- Gibbs free energy, four cases, formation route, kinetic control.
\documentclass{apchem-worksheet}

\apcunit{9}
\apcunittitle{Thermodynamics and Electrochemistry}
\apcdoctitle{Worksheet 2 \textbullet\ Gibbs Free Energy}
\apcsubtitle{CED 9.3--9.4 \textbullet\ convert the entropy to kJ before you subtract}

\begin{document}
\wsheader[$\Delta G^\circ = \Delta H^\circ - T\Delta S^\circ$ \textbullet\
crossover $T = \Delta H^\circ/\Delta S^\circ$ \textbullet\
$\Delta G^\circ = \sum n\Delta G_f^\circ{}_{\text{prod}} -
\sum n\Delta G_f^\circ{}_{\text{react}}$ \textbullet\ $\Delta H^\circ$ in
kJ, $\Delta S^\circ$ in J/K \textbullet\ take $T = \SI{298}{\kelvin}$ unless
told otherwise.]

\problem[]{\ced{9.3} The four cases. Complete the table without
calculating:}
\begin{parts}
  \item $\Delta H^\circ < 0$, $\Delta S^\circ > 0$:
        \answerblank[5.6cm]{favored at all temperatures}
  \item $\Delta H^\circ > 0$, $\Delta S^\circ < 0$:
        \answerblank[5.4cm]{favored at no temperature}
  \item $\Delta H^\circ > 0$, $\Delta S^\circ > 0$:
        \answerblank[5cm]{favored at high $T$}
  \item $\Delta H^\circ < 0$, $\Delta S^\circ < 0$:
        \answerblank[5cm]{favored at low $T$}
  \item Which two cases need a crossover temperature calculated, and why?
        \answerlines[2]{(c) and (d) --- the two terms oppose each other, so
        which one wins depends on $T$. In (a) and (b) both terms point the
        same way and temperature cannot change the outcome.}
\end{parts}

\problem[]{\ced{9.3} \ce{CaCO3(s) -> CaO(s) + CO2(g)} has
$\Delta H^\circ = +178.5$~kJ and $\Delta S^\circ = +161$~J/K.}
\begin{parts}
  \item Which of the four cases is this? \answerblank[4cm]{high $T$ only}
  \item Calculate $\Delta G^\circ$ at \SI{298}{\kelvin}. \workspace[2cm]
        \keyonly{\begin{apcanswer}$178.5 - 298(0.161) = 178.5 - 48.0 =
        \mathbf{+131}$~kJ --- not favored\end{apcanswer}}
  \item Calculate the crossover temperature. \workspace[1.8cm]
        \keyonly{\begin{apcanswer}$T = 178.5/0.161 = \mathbf{1109}$~K
        $\approx \SI{836}{\celsius}$\end{apcanswer}}
  \item Relate your answer to why limestone buildings survive but lime
        kilns work. \answerlines[2]{At ordinary temperatures
        $\Delta G^\circ$ is strongly positive, so limestone does not
        decompose. A kiln is run above about \SI{1100}{\kelvin}, where the
        $T\Delta S^\circ$ term finally exceeds $\Delta H^\circ$.}
\end{parts}

\problem[]{\ced{9.3} The Haber process has $\Delta H^\circ = -92$~kJ and
$\Delta S^\circ = -199$~J/K.}
\begin{parts}
  \item $\Delta G^\circ$ at \SI{298}{\kelvin}: \workspace[1.8cm]
        \keyonly{\begin{apcanswer}$-92 - 298(-0.199) = -92 + 59 =
        \mathbf{-33}$~kJ --- favored\end{apcanswer}}
  \item Crossover temperature: \answerblank[2.6cm]{462 K}
  \item Is the reaction favored above or below that temperature? Explain.
        \answerlines[2]{Below. Both terms are negative, so this is the
        low-temperature case: raising $T$ makes $-T\Delta S^\circ$
        increasingly positive until it overwhelms the negative
        $\Delta H^\circ$.}
  \item Industrial ammonia synthesis is run at about \SI{700}{\kelvin},
        above the crossover. Suggest why, given Unit 5.
        \answerlines[3]{At \SI{298}{\kelvin} the reaction is favored but
        far too slow. Raising the temperature sacrifices thermodynamic
        favorability for rate, and the loss in yield is recovered by high
        pressure and by recycling unreacted gas. This is a thermodynamics
        versus kinetics trade-off, not a contradiction.}
\end{parts}

\problem[]{\ced{9.3} For \ce{H2O(l) -> H2O(g)}, $\Delta H^\circ = +44$~kJ
and $\Delta S^\circ = +119$~J/K. Calculate the crossover temperature and
state what it physically is. \workspace[2.4cm]
\keyonly{\begin{apcanswer}$T = 44/0.119 = \mathbf{370}$~K
$\approx \SI{97}{\celsius}$ --- the \textbf{normal boiling point}. The
tabulated value is 373 K; the difference is rounding in the
table.\end{apcanswer}}}

\problem[]{\ced{9.3} A student calculates
$\Delta G^\circ = 178.5 - 298(161) = -47{,}800$~kJ for problem 2.
Identify the error and state how you would catch it without redoing the
arithmetic.
\answerlines[3]{They failed to convert $\Delta S^\circ$ from J/K to kJ/K,
so the entropy term is a thousand times too large. The magnitude is the
giveaway --- no ordinary reaction has a free energy change of tens of
thousands of kilojoules per mole, so a result that size signals a unit
error rather than an arithmetic one.}}

\problem[]{\ced{9.3} The formation route. For
\ce{2H2(g) + O2(g) -> 2H2O(l)}, $\Delta G_f^\circ(\ce{H2O(l)}) = -237$~kJ/mol.}
\begin{parts}
  \item Why is $\Delta G_f^\circ$ for \ce{H2(g)} and \ce{O2(g)} zero?
        \answerblank[7.2cm]{elements in their standard states}
  \item Calculate $\Delta G^\circ$. \answerblank[3cm]{$-474$ kJ}
  \item Route 1 gives $-475$~kJ from $\Delta H^\circ = -572$~kJ and
        $\Delta S^\circ = -327$~J/K. Verify, and account for the
        difference. \workspace[2.2cm]
        \keyonly{\begin{apcanswer}$-572 - 298(-0.327) = -572 + 97.4 =
        \mathbf{-475}$~kJ. The 1 kJ gap is rounding in the tabulated
        values, not an error in either method.\end{apcanswer}}
  \item State the one place $S^\circ$ and $\Delta G_f^\circ$ differ in
        convention. \answerlines[2]{$\Delta G_f^\circ$ of an element in its
        standard state is zero by definition; $S^\circ$ of that same
        element is a positive measured number.}
\end{parts}

\problem[]{\ced{9.4} Thermodynamic versus kinetic control.}
\begin{parts}
  \item $\Delta G^\circ$ tells you \answerblank[5.4cm]{whether, and how
        far}; $E_a$ tells you \answerblank[2.4cm]{how fast}.
  \item Diamond converting to graphite has $\Delta G^\circ < 0$. Explain
        why jewellery survives. \answerlines[3]{The conversion requires
        breaking and reforming a covalent network, which carries an
        enormous activation energy. The reaction is thermodynamically
        favored but under kinetic control, so its rate at room temperature
        is immeasurably small.}
  \item A catalyst changes \answerblank[2cm]{$E_a$} but does not change
        \answerblank[5.4cm]{$\Delta G^\circ$, $\Delta H^\circ$, or $K$}.
  \item Can a catalyst make an unfavorable reaction favorable? Justify.
        \answerlines[2]{No. It alters only the barrier between reactants
        and products, not their relative free energies, so
        $\Delta G^\circ$ and $K$ are untouched.}
\end{parts}

\problem[]{\ced{9.4} A reaction has $\Delta G^\circ = -200$~kJ and no
observable rate at \SI{298}{\kelvin}. A student concludes the value of
$\Delta G^\circ$ must be wrong. Explain why the conclusion does not follow,
and give two ways to make the reaction proceed.
\answerlines[4]{$\Delta G^\circ$ compares the free energies of reactants
and products and says nothing about the barrier between them, so a large
negative value is entirely compatible with a negligible rate. Raise the
temperature, so more molecules exceed $E_a$ --- this also changes $K$,
since $K$ is temperature-dependent. Or add a catalyst, which lowers $E_a$
and leaves $K$ unchanged.}}

\begin{spiralstrip}{Unit 9 \textbullet\ blocks 1--2}
  \item Sign of $\Delta S^\circ$ for \ce{2NO2(g) -> N2O4(g)}:
        \answerblank[2cm]{negative}
  \item $\Delta S^\circ$ for \ce{CaCO3(s) -> CaO(s) + CO2(g)}:
        \answerblank[2.6cm]{$+161$ J/K}
  \item $S^\circ$ of an element in its standard state:
        \answerblank[3.4cm]{not zero}
  \item Second law, as an inequality:
        \answerblank[4cm]{$\Delta S_{\text{univ}} > 0$}
  \item The AP term replacing ``spontaneous'':
        \answerblank[5.4cm]{thermodynamically favored}
\end{spiralstrip}

\end{document}
